\exercisesection{REGULAR LOCAL RINGS}

\begin{enumerate}
\def\labelenumi{\arabic{enumi})}
\item
  Let A be a ring, S a multiplicative subset of A. Show that one has \(\mathrm{dh}(\mathrm{S}^{-1}\mathrm{A}) \leqslant \mathrm{dh}(\mathrm{A})\) (A, X, p.138, def. 2). (Take resolutions of the \(S^{-1}A\) -modules by projective A-modules and localize.)
\item
  Let A be a Noetherian ring and B a faithfully flat A-algebra. Then we have \(\mathrm{dh(A)}\leqslant \mathrm{dh(B)}\).
\end{enumerate}

3) Let A be a Noetherian local ring. We propose to show that A is regular if and only if dh(A) is finite.

\begin{enumerate}
\def\labelenumi{\alph{enumi})}
\tightlist
\item
  Suppose dh(A) is finite. To show that A is regular, use A, X, p.208, exerc. 13.
\item
  Suppose A is regular. Let x be a completely secant family generating m\_A. Then, the Koszul complex associated with x is a finite resolution of A/m\_A by free A-modules. Then apply A, X, p.203, exerc. 16.
\end{enumerate}

\begin{enumerate}
\def\labelenumi{\arabic{enumi})}
\setcounter{enumi}{3}
\item
  Let A be a regular Noetherian local ring. Then we have \(\dim(A) = \mathrm{dh}(A)\) (To show that we have \(\mathrm{dh}(A) \leqslant \dim(A)\), use the Koszul complex associated with a system of coordinates. Then determine the groups \(\operatorname{Ext}_{A}^{i}(\kappa_{A}, \kappa_{A})\).)
\item
  Let A be a regular Noetherian local ring, and let p be a prime ideal of A. Then \(A_{p}\) is regular (use Exerc. 1 and Exerc. 3).
\item
  A Noetherian ring A is said to be regular if for every prime ideal p of A, the local ring \(A_{p}\) is regular.
\end{enumerate}

\begin{enumerate}
\def\labelenumi{\alph{enumi})}
\item
  Let A be a Noetherian ring. Then A is regular if and only if \(A_{m}\) is regular for every maximal ideal m of A (use Exercise 5).
\item
  Let A be a Noetherian ring. Show that A is regular of finite dimension if and only if dh(A) is finite (use Exercise 4).
\item
  Show that the ring described in Exercise 13, p.83 is regular of infinite dimension.
\item
  Let A be a regular Noetherian ring, and S a multiplicative subset. Show that \(\mathbf{S}^{-1}\mathbf{A}\) is regular.
\item
  Let A be a regular Noetherian ring. Show that A{[}X{]} is regular. (Reduce by localization to the case where dh(A) is finite. Use A, X, p. 143, th. 1, and b).
\end{enumerate}

\(^{1}\) For further details on this exercise, one may consult R. Stanley, The Hilbert function of a graded algebra, Advances in Math., 28 (1978), pp. 57–83.

\begin{enumerate}
\def\labelenumi{\arabic{enumi})}
\setcounter{enumi}{6}
\tightlist
\item
  Let I be a directed set admitting a countable cofinal subset, and let A be a ring (not necessarily commutative). We propose to show that every inductive limit, indexed by I, of projective (left) A-modules has projective dimension \(\leqslant 1\) (A, X, p.~134).
\end{enumerate}

\begin{enumerate}
\def\labelenumi{\alph{enumi})}
\tightlist
\item
  Let \(((\mathbf{Q}_i)_{i\in \mathbf{I}},(\varphi_{i,j})_{i < j})\) be an inductive system of projective A-modules and \(Q = \varinjlim_{I} Q_i\). To show that \(Q\) is of projective dimension \(\leqslant 1\), reduce to the case where \(I = \mathbf{N}\). We then have an exact sequence
\end{enumerate}

\[
0 \longrightarrow \bigoplus_ {n \in \mathbf {N}} Q _ {n} \stackrel {u} {\longrightarrow} \bigoplus_ {n \in \mathbf {N}} Q _ {n} \longrightarrow Q \longrightarrow 0
\]

where for every \(n \in \mathbf{N}\) and every \(x \in \mathbf{Q}_n\) , one has

\[
u (x) = \varphi_ {n, n + 1} (x) - x.
\]

\begin{enumerate}
\def\labelenumi{\alph{enumi})}
\setcounter{enumi}{1}
\tightlist
\item
  Let \(((\mathbf{P}_i)_{i\in I}, (\varphi_{i,j})_{i < j})\) be an inductive system of projective A-modules such that for every pair \((i,j)\) with \(i < j\), \(\varphi_{i,j}\) is injective and the module \(\operatorname{Coker}(\varphi_{i,j})\) is projective. Show that the module \(\mathbf{P} = \varinjlim_{I} \mathbf{P}_i\) is projective. (Let \(0 \to \mathbf{M}' \to \mathbf{M} \to \mathbf{M}'' \to 0\) be an exact sequence of A-modules. Then
\end{enumerate}

\[
0 \to \operatorname{Hom} _ {\mathbf {A}} (\mathrm{P} _ {i}, \mathrm{M} ^ {\prime}) \to \operatorname{Hom} _ {\mathbf {A}} (\mathrm{P} _ {i}, \mathrm{M}) \to \operatorname{Hom} _ {\mathbf {A}} (\mathrm{P} _ {i}, \mathrm{M} ^ {\prime \prime}) \to 0
\]

is an exact sequence of projective systems of commutative groups and the projective system \(i \mapsto \mathrm{Hom}_{\mathbf{A}}(\mathbf{P}_{i}, \mathbf{M}')\) has the Mittag-Leffler property for the discrete topology (TG, II, p.18). From this one will deduce that the sequence of projective limits is exact.)

\begin{enumerate}
\def\labelenumi{\alph{enumi})}
\setcounter{enumi}{2}
\tightlist
\item
  Let \((\mathbf{Q}_i)_{i\in I}\) be an inductive system of projective A-modules. For every \(i\in \mathbf{I}\), let \(\mathbf{R}_i = \bigoplus_{j\leqslant i}\mathbf{Q}_j\) and for \(i < i'\) let \(\psi_{i,i'}: \mathbf{R}_i \to \mathbf{R}_{i'}\) be the canonical injection. Define an exact sequence of inductive systems \(0 \to (\mathbf{P}_i) \to (\mathbf{R}_i) \to (\mathbf{Q}_i) \to 0\) such that \((\mathbf{P}_i)\) satisfies the hypotheses of \(b\). Deduce from this another proof of \(a\).)
\end{enumerate}

\begin{enumerate}
\def\labelenumi{\arabic{enumi})}
\setcounter{enumi}{7}
\tightlist
\item
  Let A be a ring and n an integer ≥ 1. Show that the following conditions are equivalent:
\end{enumerate}

\begin{enumerate}
\def\labelenumi{(\roman{enumi})}
\item
  We have dh(A) ≤ n.
\item
  For every cyclic A-module M, we have \(\mathrm{dp}_{\mathbf{A}}(\mathbf{M}) \leqslant n\) .
\item
  For every ideal I of A, we have \(\mathrm{dp}_{\mathbf{A}}(\mathbf{I}) \leqslant n - 1\) .
\end{enumerate}

(Use A, X, p.138, prop. 4 to prove the equivalence (i) \(\Leftrightarrow\) (ii) and A, X, p.93, prop. 11 to prove (i) \(\Leftrightarrow\) (iii).

\begin{enumerate}
\def\labelenumi{\arabic{enumi})}
\setcounter{enumi}{8}
\tightlist
\item
  Let \(\Gamma\) be a totally ordered commutative group and \(\Gamma^{+}\) the set of positive elements of \(\Gamma\). We denote by (GOD) the following property:
\end{enumerate}

(GOD): Every major subset M ⊂ Γ⁺ (VI, § 3, no. 5, def. 2) has a countable subset that is cofinal for the opposite order.

\begin{enumerate}
\def\labelenumi{\alph{enumi})}
\item
  Show that if \(\Gamma\) is of height 1, property (GOD) is satisfied.
\item
  For every integer n, show that there exists a totally ordered group \(\Gamma\) of height n possessing property (GOD) (take \(\Gamma = \mathbf{Z}^{n}\) equipped with the lexicographic order).
\end{enumerate}

10) Let A be a valuation ring whose ordered group \(\Gamma\) of values has property (GOD). a) Show that the homological dimension of A is at most 2. (Note that every ideal of A is a countable inductive limit of principal ideals, and use the results of Exercises 7 and 8.)

\begin{enumerate}
\def\labelenumi{\alph{enumi})}
\setcounter{enumi}{1}
\item
  Every valuation ring of Krull dimension 1 has homological dimension \(\leqslant 2\). It has homological dimension 2 if and only if it is not Noetherian. (Use exerc. 9, a) and A, X, p.208, exerc. 12.)
\item
  For every integer \(n \geqslant 1\), there exists a ring of Krull dimension equal to n and of homological dimension equal to 2. (Use exerc. 9, b) and A, X, p.208, exerc. 12.)
\end{enumerate}

\begin{enumerate}
\def\labelenumi{\arabic{enumi})}
\setcounter{enumi}{10}
\item
  Every regular Noetherian local ring is factorial. (Use VII, § 4, n° 7, cor. 3.)
\item
  Let A be a ring and let a be a finitely generated ideal of A.
\end{enumerate}

\begin{enumerate}
\def\labelenumi{\alph{enumi})}
\item
  Show that if there exists an ideal b of A such that a = a·b, then there exists \(b \in b\) such that \((1 - b)a = (0)\).
\item
  Show that if \(a^{2} = a\), there exists an element \(a \in a\) such that \(a^{2} = a\) and a = Aa.
\item
  We assume that a is maximal and that the A-module \(a/a^{2}\) can be generated by n elements. Show that the ideal a can be generated by \(n + 1\) elements. (One takes elements \(x_{1}, \ldots, x_{n}\) in a whose images in \(a/a^{2}\) generate this latter. Denoting by x the ideal of A generated by \(\{x_{1}, \ldots, x_{n}\}\) ; one will note that in A/x one has \((\mathfrak{a}/\mathfrak{x})^{2} = (\mathfrak{a}/\mathfrak{x})\) .)
\end{enumerate}

\begin{enumerate}
\def\labelenumi{\arabic{enumi})}
\setcounter{enumi}{12}
\tightlist
\item
  Let \(p\) be a prime number and \(f\) an integer \(> 0\). Put \(L = \mathbf{Q}(\zeta)\), where \(\zeta\) is a primitive \(p^f\)-th root of 1, and let \(B = \mathbf{Z}[\zeta]\) be the subring of L generated by \(\zeta\). Show that B is the integral closure of Z in L. (By a discriminant computation, one shows that \(B\left[\frac{1}{p}\right]\) is the integral closure of \(\mathbf{Z}\left[\frac{1}{p}\right]\) (V, § 1, no. 6, lemma 3); denoting \(\overline{\mathbf{B}}\) the integral closure of \(\mathbf{B}\),
\end{enumerate}

we deduce from this that \(\overline{\mathbf{B}}\subset\bigcap_{q}\mathbf{Z}_{q}[\zeta]\) where \(q\) runs through the prime numbers, whence \(\overline{\mathbf{B}}\subset\mathbf{Z}[\zeta].\)

\begin{enumerate}
\def\labelenumi{\arabic{enumi})}
\setcounter{enumi}{13}
\tightlist
\item
  Let \(\rho: A \to B\) be a homomorphism of Noetherian rings making \(B\) an \(A\)-module that is faithfully flat.
\end{enumerate}

\begin{enumerate}
\def\labelenumi{\alph{enumi})}
\item
  If B is regular, then A is regular (cf. p. 94, exercises 2 and 3).
\item
  If A is regular and if, for every maximal ideal m of A, B/mB is regular, then B is regular. (By localizing, reduce to the case where A and B are local and ρ is local. One then has dim(B) = dim(A) + dim(B/mB) (§ 3, no. 4, cor. 1 to prop. 7). Construct a regular sequence for B having dim(B) terms.)
\end{enumerate}

\begin{enumerate}
\def\labelenumi{\arabic{enumi})}
\setcounter{enumi}{14}
\item
  Let X be a Noetherian space and U a subset of X. If for every irreducible closed subset Y of X meeting U, U ∩ Y contains a nonempty open subset of Y, then U is open in X. (Assuming for contradiction that U is not open, consider a minimal closed set Z ⊂ X such that Z ∩ U is not open.)
\item
  Let A be a Noetherian ring. The singular locus of Spec(A), denoted Sing(A), is the set of \(p \in \text{Spec}(A)\) such that \(A_p\) is not regular. We set
\end{enumerate}

\[
\operatorname{Reg} (\mathrm{A}) = \operatorname{Spec} (\mathrm{A}) - \operatorname{Sing} (\mathrm{A}).
\]

Show that the following conditions are equivalent:

\begin{enumerate}
\def\labelenumi{(\roman{enumi})}
\item
  For every quotient ring B of A, Reg(B) is open in Spec(B).
\item
  For every integral quotient ring B of A, Reg(B) is a nonempty open subset of Spec(B).
\item
  For every integral quotient ring B of A, Reg(B) contains a nonempty open subset of Spec(B).
\end{enumerate}

(To prove (i) \(\Rightarrow\) (ii), note that one has \((0)\in\mathrm{Reg}(\mathbf{B})\) . To prove (iii) \(\Rightarrow\) (i), show that for every \(\mathfrak{p}\in\mathrm{Spec}(\mathbf{A})\) such that \(\mathrm{V}(\mathfrak{p})\cap\mathrm{Reg}(\mathbf{A})\neq\emptyset\) , the ring \(\mathbf{A}_{\mathfrak{p}}\) is regular (p.94, exerc. 5). Using (iii) and constructing completely secant sequences, show that \(\mathrm{V}(\mathfrak{p})\cap\mathrm{Reg}(\mathbf{A})\) contains a nonempty open subset of \(\mathrm{V}(\mathfrak{p})\) . Conclude using exerc. 15.)

\begin{enumerate}
\def\labelenumi{\arabic{enumi})}
\setcounter{enumi}{16}
\tightlist
\item
  Let \(\rho: A \to B\) be an injective homomorphism of integral Noetherian rings making \(B\) an \(A\)-module of finite type.
\end{enumerate}

\begin{enumerate}
\def\labelenumi{\alph{enumi})}
\item
  We assume that \(\operatorname{Reg}(\mathbf{B})\) contains a nonempty open subset of \(\operatorname{Spec}(\mathbf{B})\) . Show that \(\operatorname{Reg}(\mathbf{A})\) contains a nonempty open subset of \(\operatorname{Spec}(\mathbf{A})\) (Reduce to the case where \(\mathbf{B}\) is flat over \(\mathbf{A}\) using II, § 5, no. 1, corollary to prop. 2, then in the case where \(\mathbf{B}\) is regular by localizing with respect to a suitable element \(x \in \mathbf{A}\) Conclude using Exerc. 14, a).
\item
  We assume that \(\operatorname{Reg}(\mathbf{A})\) contains a nonempty open subset of \(\operatorname{Spec}(\mathbf{A})\) and that the field \(\mathbf{K}'\) of fractions of \(\mathbf{B}\) is a separable extension of the field \(\mathbf{K}\) of fractions of \(\mathbf{A}\) . Show that \(\operatorname{Reg}(\mathbf{B})\) contains a nonempty open subset of \(\operatorname{Spec}(\mathbf{B})\) . (Let \(\mathbf{Z} \in \mathbf{K}'\) such that \(\mathbf{K}' = \mathbf{K}(\mathbf{Z})\) (A, V, p.39), and let \(\mathrm{P}(\mathrm{X}) \in \mathrm{K}[\mathrm{X}]\) be the minimal polynomial of Z. Replacing A by the ring \(A_{f}\) for a suitable element f of A, show that one can reduce to proving the desired property of Spec(B) in the following case:
\end{enumerate}

\(\alpha)\) \(\mathbf{Z}\in \mathbf{B}\) and A is regular;

\(\beta)\) we have \(\mathrm{P}(\mathrm{X})\in \mathrm{A}[\mathrm{X}]\) ;

\(\gamma)\) we have \(\mathbf{B} = \mathbf{A}[\mathbf{X}] / (\mathbf{P}(\mathbf{X}))\) (II, § 5, no. 1, prop. 2);

δ) P'(Z) is invertible in B (observe that P'(Z) ≠ 0 in K' because K' is separable over K. Then use II, § 5, no. 1, corollary to prop. 3).

Show that in this case, for every maximal ideal m of A, B/mB is an étale (A/m)-algebra (A, V, p. 32), hence regular. Conclude using exerc. 14, b).

18) Let \(k\) be a field and A a \(k\)-algebra of finite type. Show that \(\operatorname{Reg}(A)\) is an open subset of \(\operatorname{Spec}(A)\). (Using exerc. 16, reduce to proving that \(\operatorname{Reg}(A)\) contains a nonempty open subset when A is integral. Using the normalization lemma (V, § 3, no. 1, th. 1) and exerc. 17, a), reduce to proving this last property when A is the integral closure of \(k[X_1, ..., X_n]\) in a finite quasi-Galois extension K of \(k(X_1, ..., X_n)\). Let K' be the purely inseparable closure of \(k(X_1, ..., X_n)\) in K and A' the integral closure of \(k[X_1, ..., X_n]\) in K'. Reduce to proving the desired property for A' by means of exerc. 17, b). Drawing inspiration from the proof of th. 2 of V, § 3, no. 2, embed K' in an extension \(k'(X_1^{q^{-1}}, ..., X_n^{q^{-1}}) = K''\), where \(k'\) is a purely inseparable extension of \(k\), and reduce to the case where A′ is the integral closure of \(k[X_1, ..., X_n]\) in K'' by virtue of exerc. 17, a). Then observe that A' = \(k'[X_1^{q^{-1}}, ..., X_n^{q^{-1}}]\) (V, § 1, No. 3, cor. 2 to prop. 13) and that A' is regular (p. 94, exerc. 6).

\begin{enumerate}
\def\labelenumi{\arabic{enumi})}
\setcounter{enumi}{18}
\item
  Let A be an integral domain of finite type over a field k. The set of maximal ideals \(m \in \text{Specmax}(A)\) such that \(A_{m}\) is regular is open and dense in \(\text{Specmax}(A)\).
\item
  Let \(k\) be a field, \(k_{i}\) (\(i = 1, 2\)) two extensions of \(k\), and \(n_{i}\) (\(i = 1, 2\)) the transcendence degree of \(k_{i}\) over \(k\) if the latter is finite, or the symbol +\(\infty\) otherwise. Show that we have
\end{enumerate}

\[
\dim (k _ {1} \otimes_ {k} k _ {2}) = \inf (n _ {1}, n _ {2})
\]

and that \(k_{1} \otimes_{k} k_{2}\) is Noetherian if one of the extensions \(k_{i}\) is of finite type.

\begin{enumerate}
\def\labelenumi{\arabic{enumi})}
\setcounter{enumi}{20}
\item
  Let I be a directed ordered set, \((\mathbf{A}_{i}, \varphi_{ij})\) an inductive system of rings with inductive limit A. Suppose that, for every i, the ring \(A_{i}\) is Noetherian, that for every pair \((i, j)\) such that i \textless{} j, \(\varphi_{ij}\) makes \(A_{j}\) a flat \(A_{i}\)-module, and that A is Noetherian. Show that if the \(A_{i}\) are regular for all \(i \in I\), then A is regular. (One may first reduce to the case where the \(A_{i}\) and A are local and the \(\varphi_{ij}\) are local. One will then use exercises 3, p.94 and 14, p.96.)
\item
  \begin{enumerate}
  \def\labelenumii{\alph{enumii})}
  \tightlist
  \item
    Let \(k\) be a field, A a \(k\)-algebra of finite type, and \(k'\) a finite separable extension of \(k\). Show that if A is regular, then \(A \otimes_k k'\) is regular. (One may use exercise 14, b), p.96.)
  \end{enumerate}
\end{enumerate}

\begin{enumerate}
\def\labelenumi{\alph{enumi})}
\setcounter{enumi}{1}
\tightlist
\item
  With the same hypotheses on A, let \(k'\) be a finite purely inseparable extension of \(k\). Show that in general \(A \otimes_k k'\) is not regular (take \(A = k'\)).
\end{enumerate}

\begin{enumerate}
\def\labelenumi{\arabic{enumi})}
\setcounter{enumi}{22}
\tightlist
\item
  Let \(k\) be a field of characteristic exponent \(p\), A a \(k\)-algebra of finite type, and \(\mathfrak{p}\) a prime ideal of A. Show that the following conditions are equivalent:
\end{enumerate}

\begin{enumerate}
\def\labelenumi{(\roman{enumi})}
\item
  the ring \(A_{\mathfrak{p}} \otimes_{k} k'\) is regular for every extension \(k'\) of k;
\item
  the ring \(A_{\mathfrak{p}} \otimes_{k} k^{p^{-\infty}}\) is regular;
\item
  the ring \(\mathbf{A}_{\mathfrak{p}} \otimes_{k} k'\) is regular for every finite and purely inseparable extension \(k'\) of \(k\) .
\end{enumerate}

(To prove (ii) \(\Rightarrow\) (iii), use exercise 14, a), p.96. To prove (iii) \(\Rightarrow\) (i), it suffices to show that \(A_{\mathfrak p} \otimes_k k'\) is regular when \(k'\) is a finitely generated extension of \(k\) (Exerc. 21), and when \(k'\) is a purely inseparable extension of a purely transcendental extension (Exerc. 22 and 14, a), and is bounded by a finite quasi-Galois extension of a purely transcendental extension. Then bound \(k'\) by an extension \(k''(\mathrm{X}_1, ..., \mathrm{X}_n)\), where \(k''\) is a finite purely inseparable extension of \(k\), and reduce to the case \(k' = k''(\mathrm{X}_1, ..., \mathrm{X}_n)\). By exercise 14, a), setting \(\mathbf{B} = \mathbf{A}_{\mathfrak{p}} \otimes_k k''\), we obtain a regular algebra by (iii), and \(\mathbf{A}_{\mathfrak{p}} \otimes_k k' = \mathbf{B} \otimes_{k''} k''(\mathbf{X}_1, ..., \mathbf{X}_n)\) is a regular algebra by exerc. 6, e), p.94.)

We say that \(p \in \text{Spec}(A)\) is smooth if it possesses the equivalent properties above.

\begin{enumerate}
\def\labelenumi{\arabic{enumi})}
\setcounter{enumi}{23}
\item
  Let \(k\) be a field, A a \(k\)-algebra, \(k'\) a purely inseparable extension of \(k\), and \(\mathbf{A}' = \mathbf{A} \otimes_k k'\). Show that the canonical map \(\operatorname{Spec}(\mathbf{A}') \to \operatorname{Spec}(\mathbf{A})\) is a homeomorphism. (When A is a field, \(\mathbf{A}'\) has dimension 0 (exerc. 20); one will show that it possesses only one prime ideal. From this one will deduce that for every prime ideal \(\mathfrak{p}\) of A, the radical \(\mathfrak{p}'\) of \(\mathfrak{p}\mathbf{A}'\) is a prime ideal and that the map \(\mathfrak{p} \mapsto \mathfrak{p}'\) is continuous.)
\item
  Let \(k\) be a field, A a \(k\)-algebra of finite type, and \(\operatorname{Lis}(A)\) the set of \(\mathfrak{p} \in \operatorname{Spec}(A)\) which are smooth (exerc. 23). Show that \(\operatorname{Lis}(A)\) is open in \(\operatorname{Spec}(A)\). (Use exercs. 18, 23 and 24.)
\item
  Let \(k\) be a field and A a \(k\)-algebra that is an integral domain of finite type. Show that Lis(A) is dense if and only if A is a separable \(k\)-algebra; equivalently, the field of fractions of A is a separable extension of \(k\) (A, V, p.115).
\item
  Let \(k\) be a field and A a \(k\)-algebra of finite type that is integral and separable. The set of maximal ideals of A that are smooth is an open dense subset of Specmax(A).
\item
  Let \(k\) be a field and \(k'\) an extension of \(k\). Show that for every extension \(k''\) of \(k\), \(k' \otimes_k k''\) is a regular Noetherian ring if and only if \(k'\) is a separable extension of finite type of \(k\).
\end{enumerate}

29) Let k be a field, A a complete Noetherian local k-algebra, \(\kappa_{A}\) the residue field of A.

\begin{enumerate}
\def\labelenumi{\alph{enumi})}
\item
  Suppose that \(\kappa_{\mathbf{A}}\) is a separable extension of \(k\). Show that there exists a homomorphism of \(k\)-algebras from \(\kappa_{\mathbf{A}}\) to A, which is a section of the canonical homomorphism from A to \(\kappa_{\mathbf{A}}\) (IX, § 3, n° 2, prop. 1).
\item
  Suppose that \(\kappa_{\mathbf{A}}\) is a purely inseparable extension of \(k\). Show that there does not necessarily exist a section of \(k\)-algebras from \(\kappa_{\mathbf{A}}\) into A. (Take a field \(k\) of characteristic 2, an element \(a\) of \(k\) which is not a square (for example \(k = \mathbf{F}_2(\mathrm{T})\), \(a = \mathrm{T}\)), the ideal \(\mathfrak{p}\) of \(k[\mathrm{X}]\) generated by \(\mathrm{X}^2 + a\), and set \(\mathrm{A} = \widehat{k[\mathrm{X}]_{\mathfrak{p}}}\).)
\item
  We assume A is regular and \(\kappa_{\mathbf{A}}\) is a separable extension of \(k\). Show that A is isomorphic as a \(k\)-algebra to \(\kappa_{\mathbf{A}}[[\mathrm{T}_{1}, ..., \mathrm{T}_{n}]]\), where \(n = \dim(\mathbf{A})\).
\item
  Show that this is not necessarily the case when \(\kappa_{\mathrm{A}}\) is not a separable extension of \(k\) .
\end{enumerate}

\begin{enumerate}
\def\labelenumi{\arabic{enumi})}
\setcounter{enumi}{29}
\item
  Let \(k\) be a field and A a \(k\)-algebra, local, Noetherian, and complete. One says that A is formally smooth if, for every finite extension \(k'\) of \(k\), the ring \(A \otimes_k k'\) is regular. Show that A is formally smooth if and only if, for every finite purely inseparable extension \(k'\) of \(k\), the ring \(A \otimes_k k'\) is regular (one may use exerc. 14, p.96 and imitate the proof of exerc. 23, p.97).
\item
  Let \(k\) be a field.
\end{enumerate}

\begin{enumerate}
\def\labelenumi{\alph{enumi})}
\item
  Let A be a \(k\)-algebra, Noetherian, local, complete, and regular, whose residue field \(\kappa_{\mathrm{A}}\) is a separable extension of \(k\). Show that A is formally smooth and isomorphic to \(\kappa_{\mathrm{A}}[[\mathbf{T}_{1}, ..., \mathbf{T}_{n}]]\).
\item
  Let B be a \(k\)-algebra of finite type, \(\mathfrak{p} \in \operatorname{Spec}(\mathbf{B})\) a smooth point (p.97, exerc. 23). Show that \(\hat{\mathbf{B}}_{\mathfrak{p}}\) is a formally smooth \(k\)-algebra.
\item
  Let \(\kappa\) be a finitely generated extension of \(k\) and let \(n\in\mathbf{N}\). Show that there exists a formally smooth \(k\)-algebra with residue field \(\kappa\) and of dimension \(n\).
\item
  Let A be a formally smooth k-algebra whose residue field \(\kappa_{A}\) is not a separable extension of k. Show that A is not isomorphic as a k-algebra to \(\kappa_{A}[[T_{1},...,T_{n}]]\).
\end{enumerate}

32) Let \(k\) be a field, K a finitely generated extension of \(k\), \(n\) a positive integer, and A and B two formally smooth \(k\)-algebras of dimension \(n\), such that the extensions \(\kappa_{A}\) and \(\kappa_{B}\) of \(k\) are isomorphic to K. It is proposed to show that A and B are isomorphic \(k\)-algebras.

\begin{enumerate}
\def\labelenumi{\alph{enumi})}
\item
  Reduce to the case where K is a finite purely inseparable extension of k (exercise 29, a)).
\item
  Consider the completed algebra \(C = A \hat{\otimes}_{k} B\) of \(A \otimes_{k} B\) with respect to the ideal
\end{enumerate}

\[
\left(\mathbf {A} \otimes_ {k} \mathfrak {m} _ {\mathbf {B}}\right) + \left(\mathfrak {m} _ {\mathbf {A}} \otimes_ {k} \mathbf {B}\right).
\]

Show that C is a Noetherian, local, complete, and regular \(k\)-algebra. (For this last point, one uses the canonical homomorphism \(\rho\) from A to C. We will show that it makes C a flat A-module and that \(C \otimes_{\mathbf{A}} \kappa_{\mathbf{A}}\) is isomorphic to \(\kappa_{\mathbf{A}} \otimes_{k} \mathbf{B}\). One will then use exerc. 14, p.96.)

\begin{enumerate}
\def\labelenumi{\alph{enumi})}
\setcounter{enumi}{2}
\tightlist
\item
  Show that \(\rho\) induces an isomorphism on the residue fields and an injection of \(m_{\mathbf{A}}/m_{\mathbf{A}}^{2}\) into \(m_{\mathbf{C}}/m_{\mathbf{C}}^{2}\). Then construct a retraction \(\pi\) of \(\rho\) which is a homomorphism of \(k\)-algebras. Show that \(\pi \circ \rho': B \to A\) is an isomorphism of \(k\)-algebras. (We have denoted by \(\rho'\) the canonical homomorphism from B to C.)
\end{enumerate}

